% ============================================================
% preambulo-contrato.tex
% ------------------------------------------------------------
% CONTRATO VERSIÓN 14
% ------------------------------------------------------------
% v2: LAS PRIMERAS DEL LIBRO —\gbPortadilla, \gbPortada,
% \gbCreditos— (SC-25).
% v3: \gbNombrePortada, QUE EL PREÁMBULO DEL PROYECTO USA EN
% \gbautoriaportada.
% v4: AUTORÍA EN LA APERTURA DE LA PIEZA (SC-26).
% v5: AUTORÍA EN EL SUMARIO (SC-26).
% v6: CORTE DE LÍNEA EN TÍTULOS (SC-28): \gbCorteTitulo;
% \gbCapitulo Y \gbCapituloSinNumero, QUE LA HOJA EMITE; Y
% \gbParte, \gbEncabezadoPieza Y \gbIndice RECIBEN LA VERSIÓN
% PLANA Y LA COMPUESTA DEL TÍTULO POR SEPARADO.
% v7: NOTA AL PIE EN UN TÍTULO DE SECCIÓN (SC-29): \gbNotaTitulo.
% v8: TACHADO Y SUBRAYADO (SC-30): \gbTachado Y \gbSubrayado.
% v9: NUMERACIÓN DE FIGURAS, CUADROS Y ECUACIONES POR PIEZA:
% \gbCapituloSinNumero REINICIA LOS CONTADORES (\gbReiniciarPieza)
% Y \thefigure, \thetable Y \theequation OMITEN EL PREFIJO EN UNA
% PIEZA SIN NÚMERO.
% v10: RECUADROS (SC-41): gbRecuadro Y gbRecuadroBarra, CON EL
% COLOR gbcolorbase QUE DEFINE LA ESTÉTICA.
% v11: CÓDIGO (SC-42): EL ENTORNO gbCodigo, LAS MACROS DE LÍNEA
% \gbnl, \gbnc Y \gbret, LAS DE TOKEN DE skylighting
% (\KeywordTok Y LAS DEMÁS) Y EL PIE \gbPieCodigo CON SU
% CONTADOR gbcodigo, NUMERADO COMO LAS FIGURAS. listings SE
% RETIRA: LO REEMPLAZA fancyvrb.
% v12: CONVERSACIÓN (SC-43): EL ENTORNO gbconversacion Y LA MACRO
% \gbetiqueta. SE RETIRAN gbparlamento Y \gblocutor.
% v13: DEDICATORIA: EL ENTORNO gbdedicatoria Y LA MEDIDA
% \gbBajadaDedicatoria, QUE DEFINE LA ESTÉTICA.
% v14: VERSO (SC-49): EL PAQUETE verse; gbverso PASA A SER EL
% verse DEL PAQUETE (CADA VERSO TERMINA EN \\, CADA ESTROFA EN
% \\!, LA SANGRÍA ES \vin) Y SE AGREGA gbversoepigrafe.
% v15: CIFRAS Y CANTIDADES (SC-64, SC-65): siunitx CON SU
% \sisetup, Y LA UNIDAD \euro. EL XSLT EMITE \num, \SI Y \SIrange.
% v16: CUADROS (SC-66): \gbCuadroLetra, \gbCuadroCuerpoNormal,
% \gbCuadroCuerpoChico, \gbCuadroSepNormal Y \gbCuadroSepEstrecha.
% v17: CUADROS, ENTREGA 2 (SC-66): EL PAQUETE rotating (sidewaystable)
% Y LOS TEXTOS DE CONTINUACIÓN \gbCuadroContinua Y \gbCuadroViene.
% v18: ÍNDICES ONOMÁSTICO Y DE CONCEPTOS (SC-77): EL XSLT EMITE
% \index CON LOS FORMATOS \gbIndNota Y \gbIndPrincipal; EL ÍNDICE
% names AGREGA |hyperpage (hyperref VA CON hyperindex=false, GV-159).
% v19: SIGLAS (1.18.0): \gbextranjero, EL DESARROLLO DE UNA SIGLA EN
% OTRO IDIOMA, Y LA LISTA DE SIGLAS: \gbSiglasInicio, \gbSiglaMedir,
% EL ENTORNO gbsiglas Y LA MEDIDA \gbSiglasSep.
% v20: GLOSARIO (1.18.0): EL ENTORNO gbglosario, \gbGlosarioLetra,
% \gbGlosarioAncla, \gbGlosarioTermino, \gbGlosarioSep,
% \gbGlosarioVease, \gbGlosarioRemite Y LA MEDIDA \gbGlosarioSangria.
% v21: SE PARTE (07/11/2026). LO QUE LAS DOS HOJAS DE LaTeX
% EMITEN Y NO DEPENDE DEL CAPÍTULO PASA A preambulo-contrato-comun.tex,
% QUE ESTE ARCHIVO CARGA PRIMERO; ACÁ QUEDA LO DE LIBROS. SIN CAMBIOS
% DE COMPORTAMIENTO: gbcodigo SE DECLARA EN EL COMÚN Y ACÁ SE LE SUMA
% EL REINICIO POR CAPÍTULO, Y citereset=chapter SALE DE LOS
% cita-*-config.tex Y SE FIJA ACÁ DESPUÉS DE CARGAR biblatex.
% ------------------------------------------------------------
% DEFINE TODO LO QUE docbook-to-latex.xsl PUEDE EMITIR EN LOS
% FRAGMENTOS DE CAPÍTULO. ES LA CAPA QUE UNA EDITORIAL
% CONSERVA CUANDO REEMPLAZA LA ESTÉTICA: SI FALTA UNA DE
% ESTAS DEFINICIONES, EL LIBRO NO COMPILA; SI ESTÁ DEFINIDA
% DISTINTO, COMPILA Y SALE MAL.
%
% CUANDO SE AGREGUE VOCABULARIO NUEVO, SUBIR LA VERSIÓN DEL
% CONTRATO. EL main.tex DECLARA LA QUE NECESITA Y ESTA CAPA
% LA VERIFICA: ASÍ UN PREÁMBULO VIEJO FRENA EN LA PRIMERA
% LÍNEA CON UN MENSAJE ÚTIL, EN VEZ DE DAR UN
% "Undefined control sequence" EN LA PÁGINA 140.
% ------------------------------------------------------------
% DEPENDE DE preambulo-proyecto.tex, QUE VA ANTES.
% ============================================================

\def\gbContratoVersion{21}

% ============================================================
% \gbRequiereContrato{N}
% PROPÓSITO : EL main.tex DECLARA QUÉ VERSIÓN NECESITA.
% PARÁMETRO : N — versión requerida
% ============================================================
\newcommand{\gbRequiereContrato}[1]{%
  \ifnum#1>\gbContratoVersion\relax
    \PackageError{gbpublisher}{%
      Este preambulo implementa el contrato \gbContratoVersion\space
      y el libro requiere el #1}{%
      Actualizar preambulo-contrato.tex.}%
  \fi}

% ============================================================
% LO COMÚN CON REVISTAS (v21)
% ------------------------------------------------------------
% CORTE Y NOTA EN TÍTULOS, TACHADO Y SUBRAYADO, PAQUETES, babel,
% csquotes, CIFRAS, CUADROS, BLOQUES, LLAMADAS DE ATENCIÓN Y
% CÓDIGO. VA PRIMERO: LO DE LIBROS QUE SIGUE SON DEFINICIONES QUE
% USAN ESOS PAQUETES, Y biblatex PIDE babel Y csquotes ANTES.
% ============================================================
% ============================================================
% preambulo-contrato-comun.tex
% ------------------------------------------------------------
% CONTRATO COMÚN VERSIÓN 1
% ------------------------------------------------------------
% v1 (07/11/2026): SALE DE preambulo-contrato.tex, QUE ERA DE
% LIBROS. ES LO QUE EMITEN LAS DOS HOJAS DE LaTeX,
% docbook-to-latex.xsl (LIBROS) Y jats-to-latex.xsl (REVISTAS), Y
% NO DEPENDE DE LA CLASE NI DEL CAPÍTULO: CORTE Y NOTA EN
% TÍTULOS, TACHADO Y SUBRAYADO, LOS PAQUETES DEL CONTRATO, babel,
% csquotes, CIFRAS, CUADROS, BLOQUES (CITA, EPÍGRAFE, RECUADROS,
% VERSO, CONVERSACIÓN), LLAMADAS DE ATENCIÓN Y CÓDIGO. EL TEXTO
% DE CADA APARTADO SE MUDÓ SIN CAMBIOS, SALVO EL CONTADOR
% gbcodigo (VER «8. CÓDIGO»).
% LO CARGA, PRIMERO, EL CONTRATO DE CADA TIPO DE PROYECTO:
%   libro    preambulo-contrato.tex
%   revista  preambulo-revista-contrato.tex
% UN CAMBIO ACÁ LLEGA A LOS DOS: SE PRUEBA CON UN LIBRO Y CON
% UNA REVISTA ANTES DE SUBIR LA VERSIÓN.
% ------------------------------------------------------------
% DEPENDE DE LAS DEFINICIONES DEL PROYECTO, QUE VAN ANTES:
%   \gbidiomas          LOS IDIOMAS DE babel
%   \ifgbcodigocolor    EL CÓDIGO EN COLOR O EN GRISES (SI NO ESTÁ
%                       DEFINIDO, EN COLOR)
% ============================================================

\def\gbContratoComunVersion{1}

% ============================================================
% CORTE DE LÍNEA EN TÍTULOS (CONTRATO 6, SC-28)
% ------------------------------------------------------------
% El editor marca en la base dónde cortar un título o un
% subtítulo. El generador lo traduce a \gbCorteTitulo SOLO en
% la versión compuesta del título: la que va a la apertura de
% la pieza o de la parte y a la portada. El sumario, el folio y
% los marcadores reciben siempre la versión plana, sin la
% macro: por eso las macros de título de este contrato llevan
% las dos.
%
% \gbtitulo Y \gbsubtitulo VAN COMPUESTOS Y TAMBIÉN LLEGAN A
% pdftitle Y pdfsubject. AHÍ hyperref PASA EL VALOR A CADENA
% PDF, Y UN \\ NO TIENE LUGAR: LA DESACTIVACIÓN LO CAMBIA POR
% UN ESPACIO. VALE ESPACIO Y NO NADA PORQUE EL GENERADOR ESCRIBE
% «Uno\gbCorteTitulo dos», Y EL ESPACIO QUE SIGUE A LA MACRO LO
% CONSUME TeX AL LEER EL NOMBRE.
%
% hyperref SE CARGA EN LA ÚLTIMA CAPA, Y \pdfstringdefDisableCommands
% NO EXISTE ANTES: SE REGISTRA EN EL GANCHO DE CARGA DEL PAQUETE,
% QUE CORRE APENAS hyperref TERMINA DE CARGARSE Y ANTES DEL
% \hypersetup{pdftitle=...} DE ESA MISMA CAPA. EN \AtBeginDocument
% SERÍA TARDE: pdftitle YA ESTARÍA CONVERTIDO.
% ============================================================
\newcommand{\gbCorteTitulo}{\\}
\AddToHook{package/hyperref/after}{%
  \pdfstringdefDisableCommands{\def\gbCorteTitulo{ }}}

% ============================================================
% NOTA AL PIE EN UN TÍTULO DE SECCIÓN (CONTRATO 7, SC-29)
% ------------------------------------------------------------
% \gbNotaTitulo{TEXTO}: LA EMITE docbook-to-latex.xsl EN EL
% ARGUMENTO OBLIGATORIO DE \section Y SUS NIVELES, NUNCA EN EL
% OPCIONAL: EL SUMARIO, EL FOLIO Y LOS MARCADORES VAN SIN NOTA.
%
% CON LA OPCIÓN calcwidth, titlesec COMPONE EL TÍTULO DOS VECES,
% UNA PARA MEDIRLO, Y UN \footnote DIRECTO AVANZA DOS VECES EL
% CONTADOR: LAS NOTAS SALEN 2, 3, 5 (VERIFICADO, GV-70). EL
% MANUAL DE titlesec (§2.9) DA \iftitlemeasuring{MIDE}{COMPONE}
% PARA ESTE CASO: AL MEDIR NO SE EMITE NADA.
%
% SI LA ESTÉTICA NO CARGA titlesec, \iftitlemeasuring NO EXISTE:
% SE PROVEE AL EMPEZAR EL DOCUMENTO, CUANDO LA ESTÉTICA YA SE
% CARGÓ, Y SIEMPRE COMPONE. \providecommand NO TOCA LA DE
% titlesec.
%
% LOS PÁRRAFOS DE LA NOTA LLEGAN SEPARADOS CON \endgraf: EL
% ARGUMENTO DE UN TÍTULO NO ADMITE \par.
% ============================================================
\newcommand{\gbNotaTitulo}[1]{%
  \iftitlemeasuring{}{\footnote{#1}}}
\AtBeginDocument{%
  \providecommand{\iftitlemeasuring}[2]{#2}}

% ============================================================
% TACHADO Y SUBRAYADO (CONTRATO 8, SC-30)
% ------------------------------------------------------------
% \gbTachado{TEXTO} Y \gbSubrayado{TEXTO}: LOS EMITE
% docbook-to-latex.xsl PARA ~~x~~ Y [x]{.underline}.
%
% ulem Y NO \underline: \underline ARMA UNA CAJA Y NO CORTA LA
% LÍNEA, Y LaTeX NO TIENE TACHADO PROPIO. [normalem] ES
% OBLIGATORIO: SIN ÉL, ulem CAMBIA \emph POR SUBRAYADO EN TODO
% EL LIBRO.
%
% \DeclareRobustCommand PORQUE PUEDEN QUEDAR EN UN TÍTULO DE
% SECCIÓN, QUE VIAJA AL SUMARIO, Y LOS COMANDOS DE ulem SON
% FRÁGILES. EN LOS MARCADORES DEL PDF SE QUEDA EL TEXTO SOLO.
% ============================================================
\usepackage[normalem]{ulem}
\DeclareRobustCommand{\gbTachado}[1]{\sout{#1}}
\DeclareRobustCommand{\gbSubrayado}[1]{\uline{#1}}
\AddToHook{package/hyperref/after}{%
  \pdfstringdefDisableCommands{%
    \def\gbTachado#1{#1}%
    \def\gbSubrayado#1{#1}}}

% ============================================================
% 1. PAQUETES QUE EL CONTRATO NECESITA
% ------------------------------------------------------------
% SOLO LO QUE HACE FALTA PARA QUE UN FRAGMENTO COMPILE. LA
% ESTÉTICA CARGA LO SUYO APARTE. TODOS ESTÁN EN TeX Live Y
% NINGUNO PIDE --shell-escape.
% ============================================================
\usepackage{etoolbox}
\usepackage{graphicx}
\usepackage{booktabs}
\usepackage{array}
% fancyvrb: EL Verbatim CON commandchars DEL CÓDIGO COLOREADO (SC-42)
\usepackage{fancyvrb}
\usepackage{ragged2e}
\usepackage{xcolor}
\usepackage[most]{tcolorbox}
% SEPARADOR ORNAMENTAL \froufrou (SC-36). CARGA tikz Y fourier-orns
\usepackage{froufrou}
% VERSO (SC-49): REDEFINE verse Y TRAE \vin, \\! Y \stanzaskip. EN
% TeX Live 2023 (MINT) ES LA 2.4a; LO QUE SE USA EXISTE DESDE LA 2.2
\usepackage{verse}

% ------------------------------------------------------------
% IDIOMA
% ------------------------------------------------------------
% babel VA EN EL CONTRATO Y NO EN LA ESTÉTICA, POR DOS RAZONES.
% LA PRIMERA ES DE ORDEN: biblatex PIDE CARGARSE DESPUÉS DE
% babel Y csquotes, Y LA ESTÉTICA VA DESPUÉS DE ESTA CAPA.
% LA SEGUNDA ES DE FONDO: EL IDIOMA DECIDE LAS COMILLAS QUE
% PRODUCE \enquote, LA PARTICIÓN DE PALABRAS Y LAS CADENAS DE
% LA BIBLIOGRAFÍA. SON GARANTÍAS DEL CONTRATO, NO DECISIONES
% DE DISEÑO: UNA EDITORIAL QUE REEMPLAZA LA ESTÉTICA NO DEBE
% PODER CAMBIAR EL IDIOMA DEL LIBRO.
\usepackage[\gbidiomas]{babel}
\frenchspacing

% COMILLAS SEGÚN EL IDIOMA ACTIVO: EN ESPAÑOL DA « ».
% EL XSLT EMITE \enquote Y NUNCA COMILLAS LITERALES.
\usepackage[autostyle=true]{csquotes}

% ------------------------------------------------------------
% CIFRAS Y CANTIDADES (CONTRATO 15, SC-64, SC-65)
% ------------------------------------------------------------
% LA NORMA INTERNACIONAL (SI Brochure §5.4): PUNTO DECIMAL,
% GRUPOS DE A TRES CON UN ESPACIO DESDE CINCO DÍGITOS, LA UNIDAD
% Y EL PORCENTAJE SEPARADOS DE LA CIFRA, EL GRADO PEGADO.
% detect-all NO SE QUITA: SIN ELLA siunitx COMPONE LA CIFRA EN
% MODO MATEMÁTICO, DONDE babel-spanish PASA EL PUNTO A COMA, Y
% LOS DECIMALES SALEN «12,5» SIN AVISO (GV-123). EN siunitx 3 ES
% UNA OPCIÓN DEPRECADA: DEJA UN AVISO INFORMATIVO EN EL LOG.
% per-mode=symbol: «48 km/h» EN UNA LÍNEA; LA FRACCIÓN VERTICAL
% AGRANDABA LA LÍNEA (DECISIÓN DE ALBERTO).
% EL FORMATEADOR DE cifras.lua APLICA LAS MISMAS REGLAS PARA EL
% HTML Y EL EPUB: SI CAMBIA ESTE \sisetup, CAMBIA EL FILTRO.
\usepackage{siunitx}
\sisetup{output-decimal-marker={.},
  group-separator={\hspace{0.15em}},
  group-minimum-digits=5,
  table-text-alignment=center,
  detect-all,
  per-mode=symbol}
% siunitx NO DEFINE EL EURO (GV-123). LA FUENTE TIENE QUE TENER €
\DeclareSIUnit\euro{€}

% ------------------------------------------------------------
% CUADROS (CONTRATO 16, SC-66)
% ------------------------------------------------------------
% EL ESTILO DE LA CASA: LETRA SIN SERIFA, FILETE ENTRE TODAS LAS
% FILAS Y ENTRE TODAS LAS COLUMNAS, ENCABEZADO EN NEGRITA, PIE
% DEBAJO. LOS FILETES Y LA NEGRITA LOS ESCRIBE EL XSLT; ACÁ, LO
% QUE UNA EDITORIAL PUEDE CAMBIAR SIN TOCARLO. EL CUERPO Y LA
% SEPARACIÓN DE COLUMNAS LOS ELIGE EL MAQUETADOR POR CUADRO
% (cuerpo: normal|chico, separacion: normal|estrecha). LA LETRA
% Y EL CUERPO VALEN PARA TODO EL FLOTANTE, EL PIE INCLUIDO.
% GEMELAS: m_XML.ObtenerPreambuloEmbebido (REVISTAS).
\newcommand{\gbCuadroLetra}{\sffamily}
\newcommand{\gbCuadroCuerpoNormal}{\footnotesize}
\newcommand{\gbCuadroCuerpoChico}{\scriptsize}
\newcommand{\gbCuadroSepNormal}{4pt}
\newcommand{\gbCuadroSepEstrecha}{2pt}
% UN CUADRO CORTADO (cortes:) VA EN UN FLOTANTE POR TRAMO: EL PRIMERO
% TERMINA CON \gbCuadroContinua, A LA DERECHA; LOS QUE SIGUEN EMPIEZAN
% CON \gbCuadroViene, A LA IZQUIERDA, Y REPITEN EL ENCABEZADO. EL PIE,
% SOLO EN EL ÚLTIMO. EL XSLT LOS PONE EN BASTARDILLA.
\newcommand{\gbCuadroContinua}{Continúa en la página siguiente}
\newcommand{\gbCuadroViene}{Viene de la página anterior}
% EL CUADRO APAISADO (rotar: sí): sidewaystable, UNA PÁGINA POR TRAMO
\usepackage{rotating}

% ============================================================
% 6. BLOQUES
% ============================================================

% CITA EN BLOQUE CON ATRIBUCIÓN OPCIONAL.
% NO ES UN EPÍGRAFE: DOCBOOK LOS DISTINGUE Y EL XSLT TAMBIÉN.
\newenvironment{gbcita}
  {\begin{quote}\small}
  {\end{quote}}
\newcommand{\gbatribucion}[1]{\par\vspace{2pt}\hfill\textup{\small#1}\par}

% EPÍGRAFE (SC-35). BLOQUE A LA DERECHA, 60 % DE LA COLUMNA, \small,
% SIN CORTE DE PALABRA. #1 A LA IZQUIERDA; FILETE DE 0,6 PT; #2 A LA
% DERECHA. CON #2 VACÍO NO HAY FILETE. ENTRE PÁRRAFOS DE #1, MEDIA LÍNEA. #1 PUEDE TENER VARIOS PÁRRAFOS
% (\newcommand SIN ASTERISCO ADMITE \par). #2 PUEDE SER UNA CITA: EL
% XSLT RESUELVE UN <biblioref> A \textcite.
% GEMELA DE LA DEFINICIÓN DE m_XML.ObtenerPreambuloEmbebido (REVISTAS):
% UN CAMBIO EN UNA VA EN LAS DOS.
\newcommand{\gbepigrafe}[2]{%
  \par\vspace{.5\baselineskip}%
  {\noindent\hfill
   \begin{minipage}{.6\linewidth}%
     \small\parindent=0pt\parskip=.5em
     \hyphenpenalty=10000\exhyphenpenalty=10000
     \raggedright #1\par
     \if\relax\detokenize{#2}\relax\else
       \parskip=0pt\vspace{1ex}\hrule height 0.6pt\vspace{1ex}%
       {\raggedleft #2\par}%
     \fi
   \end{minipage}\par}%
  \vspace{.5\baselineskip}}

% RECUADRO LATERAL. UN ARGUMENTO: info | supplementary.
\newenvironment{gbsidebar}[1]
  {\begin{tcolorbox}[colback=black!4,colframe=black!4,
     boxrule=0pt,arc=2pt,left=6pt,right=6pt,top=4pt,bottom=4pt,
     before skip=6pt,after skip=6pt]\small\sffamily}
  {\end{tcolorbox}}

% ============================================================
% RECUADROS (CONTRATO 10, SC-41)
% ------------------------------------------------------------
% gbRecuadro: FONDO GRIS AL 7 %, FILETE DE 1 PT Y ESQUINAS DE
% 5 PT, LOS DOS EN gbcolorbase.
% gbRecuadroBarra{BARRA}: FONDO GRIS AL 7 %, SIN FILETE Y CON
% ESQUINAS RECTAS; ARRIBA, UNA BARRA DE UNA PICA EN gbcolorbase
% CON EL TEXTO EN SANS NEGRITA Y MAYÚSCULAS, BLANCO Y CENTRADO
% EN ALTURA: EL MISMO RÓTULO QUE EL TIPO DE ARTÍCULO DE LA
% CABECERA. \MakeUppercase Y NO upper-case() EN EL XSLT: CONSERVA
% LA BASTARDILLA Y DEJA INTACTA LA CLAVE DE UN \autocite (MEDIDO).
% EL TEXTO ES EL DEL CUERPO: breakable PARA QUE UN RECUADRO
% LARGO CRUCE DE PÁGINA; parbox=false PARA QUE LOS PÁRRAFOS
% TENGAN LA SANGRÍA DEL CUERPO (EN MODO MINIPÁGINA QUEDARÍAN
% PEGADOS Y SIN SANGRÍA); EL PRIMERO VA SIN SANGRÍA.
% SIN NOTAS AL PIE: EN UN tcolorbox QUEDAN ADENTRO, CON LETRA
% (MEDIDO); EL FILTRO recuadros.lua LAS FRENA.
% EL COLOR LO DEFINE LA ESTÉTICA (\definecolor{gbcolorbase});
% ESTE \providecolor ES EL RESPALDO DE UNA ESTÉTICA QUE NO LO
% DECLARE. GEMELOS EN m_XML.ObtenerPreambuloEmbebido (REVISTAS):
% UN CAMBIO EN UNO VA EN LOS DOS.
% ============================================================
\providecolor{gbcolorbase}{HTML}{2D5A8E}
\tcbset{
  gbrecuadro/.style={breakable, parbox=false,
    colback=black!7, colframe=gbcolorbase, boxrule=1pt, arc=5pt,
    boxsep=0pt, left=8pt, right=8pt, top=8pt, bottom=8pt,
    before skip=\baselineskip, after skip=\baselineskip,
    before upper=\noindent},
  gbrecuadrobarra/.style={breakable, parbox=false, sharp corners,
    colback=black!7, colframe=black!7, boxrule=0pt,
    colbacktitle=gbcolorbase, coltitle=white, titlerule=0pt,
    boxsep=0pt, left=8pt, right=8pt, top=8pt, bottom=8pt,
    toptitle=0pt, bottomtitle=0pt, lefttitle=8pt, righttitle=8pt,
    before skip=\baselineskip, after skip=\baselineskip,
    before upper=\noindent}}
\newenvironment{gbRecuadro}
  {\begin{tcolorbox}[gbrecuadro]}
  {\end{tcolorbox}}
\newenvironment{gbRecuadroBarra}[1]
  {\begin{tcolorbox}[gbrecuadrobarra,
     title={\parbox[c][1pc][c]{\linewidth}{\sffamily\bfseries\MakeUppercase{#1}}}]}
  {\end{tcolorbox}}

% ============================================================
% VERSO (CONTRATO 14, SC-49)
% ------------------------------------------------------------
% EL ENTORNO verse DEL PAQUETE verse. EL XSLT TERMINA CADA VERSO
% EN \\ Y EL ÚLTIMO DE CADA ESTROFA EN \\!, QUE DEJA \stanzaskip
% (0,75 DE LÍNEA) ENTRE ESTROFAS; EL ÚLTIMO DEL POEMA VA SIN NADA.
% LA SANGRÍA ES \vin (\vgap, 1,5 em) UNA VEZ POR NIVEL; UN VERSO
% QUE NO ENTRA SIGUE A \vindent (3 em). CADA VERSO EMPIEZA CON
% \relax, PARA QUE UN [ O UN * DEL TEXTO NO SE LEA COMO ARGUMENTO
% DE \\ (GV-93).
% gbverso: EL POEMA, EN BASTARDILLA, A \leftmargini DEL MARGEN.
% gbversoepigrafe: EL VERSO DE UN EPÍGRAFE, DENTRO DEL PRIMER
% ARGUMENTO DE \gbepigrafe: LA LETRA DEL EPÍGRAFE, SIN MARGEN
% PROPIO Y SIN EL \topsep DE LA LISTA, QUE SEPARARÍA EL ÚLTIMO
% VERSO DEL FILETE. LOS DOS SALEN DE \@listi, QUE \list LEE AL
% ABRIR: SE REDEFINE LOCAL AL ENTORNO. MEDIDO: DEL CUERPO AL
% PRIMER VERSO Y DEL ÚLTIMO VERSO A LA ATRIBUCIÓN, LA MISMA
% DISTANCIA QUE EN UN EPÍGRAFE EN PROSA (GV-94). \parsep E
% \itemsep LOS PONE verse DESPUÉS (\stanzaskip ENTRE ESTROFAS).
% GEMELOS EN m_XML.ObtenerPreambuloEmbebido (REVISTAS): UN CAMBIO
% EN UNOS VA EN LOS OTROS.
% ============================================================
\newenvironment{gbverso}
  {\begin{verse}\itshape}
  {\end{verse}}
\makeatletter
\newenvironment{gbversoepigrafe}
  {\def\@listi{\leftmargin\z@\topsep\z@\partopsep\z@
     \itemsep\z@\parsep\z@}%
   \begin{verse}}
  {\end{verse}}
\makeatother

% ============================================================
% CONVERSACIÓN (CONTRATO 12, SC-43)
% ------------------------------------------------------------
% gbconversacion: EL qandaset role="conversacion" ENTERO. SANS,
% \small; TODO EL BLOQUE CORRIDO 14 PT A LA IZQUIERDA (\leftskip), EL
% MISMO MARGEN QUE LA CITA (quote, EN preambulo-estetica.tex), SIN
% MARGEN A LA DERECHA; SIN SANGRÍA DE PRIMERA LÍNEA. \leftskip Y NO
% UNA LISTA COMO LA CITA: DENTRO DE UNA LISTA EL ESPACIO ENTRE PÁRRAFOS
% ES \parsep, NO \parskip. LAS NOTAS NO LO HEREDAN: \@parboxrestore
% LO PONE EN CERO. UN ESPACIO CHICO (\smallskipamount) ENTRE
% TURNOS Y ENTRE LOS PÁRRAFOS DE UN TURNO, QUE SIN SANGRÍA ES LO
% ÚNICO QUE LOS SEPARA; \medskip AL ABRIR Y AL CERRAR, QUE DA LA
% SENSACIÓN DE BLOQUE.
% \gbetiqueta{ETIQUETA}: ABRE EL PRIMER PÁRRAFO DEL TURNO, EN
% LÍNEA. LA HOJA LA ESCRIBE YA EN MAYÚSCULAS; PESO NORMAL.
% LA PREGUNTA NO QUEDA SOLA AL PIE: LA HOJA ESCRIBE \nopagebreak[4]
% DESPUÉS DE CADA UNA.
% GEMELO DE LA DEFINICIÓN DE m_XML.ObtenerPreambuloEmbebido
% (REVISTAS): UN CAMBIO EN UNA VA EN LAS DOS.
% ============================================================
\newenvironment{gbconversacion}
  {\par\medskip
   \begingroup\sffamily\small
   \setlength{\leftskip}{14pt}%
   \setlength{\parindent}{0pt}%
   \setlength{\parskip}{\smallskipamount}}
  {\par\endgroup\medskip}
\newcommand{\gbetiqueta}[1]{#1\enspace}

% ============================================================
% 7. LLAMADAS DE ATENCIÓN
% ============================================================
% CADA LLAMADA DE ATENCIÓN ES UN tcolorbox CON TÍTULO FIJO.
% NO SE USA UN .style DE pgfkeys PARAMETRIZADO: EL #1 DE UN
% \newcommand DENTRO DE UN .style LO CONSUME pgfkeys Y DA
% "You already have nine parameters", QUE NO DICE NADA SOBRE
% LA CAUSA. CON UN ENTORNO POR CLASE EL PROBLEMA NO EXISTE.
\tcbset{gbadmon/.style={
  colback=black!3, colframe=black!25,
  boxrule=0.4pt, arc=1pt,
  left=6pt, right=6pt, top=4pt, bottom=4pt,
  before skip=6pt, after skip=6pt,
  fonttitle=\bfseries\sffamily\small}}

\newenvironment{gbnota}
  {\begin{tcolorbox}[gbadmon,title={Nota}]\small}{\end{tcolorbox}}
\newenvironment{gbconsejo}
  {\begin{tcolorbox}[gbadmon,title={Sugerencia}]\small}{\end{tcolorbox}}
\newenvironment{gbadvertencia}
  {\begin{tcolorbox}[gbadmon,title={Advertencia}]\small}{\end{tcolorbox}}
\newenvironment{gbimportante}
  {\begin{tcolorbox}[gbadmon,title={Importante}]\small}{\end{tcolorbox}}
\newenvironment{gbprecaucion}
  {\begin{tcolorbox}[gbadmon,title={Precaución}]\small}{\end{tcolorbox}}

% ============================================================
% 8. CÓDIGO (CONTRATO 11, SC-42)
% ------------------------------------------------------------
% EL COLOR NO LO PONE LaTeX: LO PONE EL RESALTADOR DE PANDOC
% (skylighting) EN UN PASO PREVIO (colorear_codigo.sh), QUE
% ESCRIBE CADA LÍNEA CON LAS MACROS DE TOKEN (\KeywordTok{for}
% Y LAS DEMÁS). ACÁ SOLO SE DEFINE CÓMO SE VE CADA UNA. NI
% listings NI minted: listings NO CONOCE JavaScript, JSON, YAML
% NI Lua (MEDIDO), Y minted EXIGE --shell-escape Y Pygments.
%
% gbCodigo{RÓTULO}: CAJA OSCURA CON CABECERA, EL RÓTULO DEL
% LENGUAJE ARRIBA Y LAS LÍNEAS EN UN Verbatim CON commandchars.
% breakable: UN BLOQUE LARGO CRUZA DE PÁGINA.
% \gbnl{N}: NÚMERO DE LÍNEA, AL MARGEN Y EN GRIS.
% \gbnc: CONTINUACIÓN DE UN CORTE MANUAL, SIN NÚMERO.
% \gbret: EL ↩ DEL CORTE MANUAL, EN EL GRIS DE LOS NÚMEROS.
% \gbPieCodigo{ID}{TEXTO}: EL PIE DE UN LISTADO, DEBAJO DEL
% BLOQUE COMO EN LAS FIGURAS: «Código 2.3.», CON EL CONTADOR
% gbcodigo (DECLARADO CON LOS DE LA NUMERACIÓN POR PIEZA).
%
% \ifgbcodigocolor: VERDADERO, PALETA Gruvbox OSCURA (EL TEMA
% DE LA APLICACIÓN); FALSO, FONDO BLANCO, TEXTO NEGRO Y SIN
% COLOR, PARA IMPRENTA. LA BANDERA LA ESCRIBE EL GENERADOR EN
% LAS PRESTACIONES DE SALIDA (estado «publicado»); SI NO LLEGA,
% VA EL COLOR. LA NEGRITA Y LA BASTARDILLA DE LOS TOKENS SE
% CONSERVAN EN LOS DOS CASOS.
%
% LA PALETA ES LA DE .hidden/themes/gruvbox.theme. GEMELOS: EL
% CSS DEL CÓDIGO EN gbpublisher.css Y EN jats-to-html.xsl, Y
% ESTA MISMA SECCIÓN EN m_XML.ObtenerPreambuloEmbebido
% (REVISTAS). UN CAMBIO DE COLOR VA EN LOS CUATRO.
% ============================================================
\expandafter\ifx\csname ifgbcodigocolor\endcsname\relax
  \expandafter\newif\csname ifgbcodigocolor\endcsname
  \gbcodigocolortrue
\fi

\ifgbcodigocolor
  \definecolor{gbcodfondo}{HTML}{282828}
  \definecolor{gbcodmarco}{HTML}{282828}
  \definecolor{gbcodtexto}{HTML}{EBDBB2}
  \definecolor{gbcodcabeza}{HTML}{3C3836}
  \definecolor{gbcodrotulo}{HTML}{A89984}
  \definecolor{gbcodnumero}{HTML}{7C6F64}
\else
  \definecolor{gbcodfondo}{gray}{1}
  \definecolor{gbcodmarco}{gray}{0}
  \definecolor{gbcodtexto}{gray}{0}
  \definecolor{gbcodcabeza}{gray}{0.9}
  \definecolor{gbcodrotulo}{gray}{0}
  \definecolor{gbcodnumero}{gray}{0.45}
\fi

% \gbDefTok{NOMBRE}{COLOR HTML}{ESTILO}: DEFINE \NOMBRETok. EL
% ESTILO ES \textbf, \textit O \relax. SIN COLOR, SOLO EL ESTILO.
\newcommand{\gbDefTok}[3]{%
  \ifgbcodigocolor
    \expandafter\newcommand\csname #1Tok\endcsname[1]{\textcolor[HTML]{#2}{#3{##1}}}%
  \else
    \expandafter\newcommand\csname #1Tok\endcsname[1]{#3{##1}}%
  \fi}
% LAS 31 CLASES DE skylighting
\gbDefTok{Normal}{EBDBB2}{\relax}
\gbDefTok{Variable}{EBDBB2}{\relax}
\gbDefTok{Keyword}{83A598}{\textbf}
\gbDefTok{ControlFlow}{83A598}{\textbf}
\gbDefTok{DataType}{FABD2F}{\relax}
\gbDefTok{Attribute}{FABD2F}{\relax}
\gbDefTok{DecVal}{D3869B}{\relax}
\gbDefTok{BaseN}{D3869B}{\relax}
\gbDefTok{Float}{D3869B}{\relax}
\gbDefTok{Constant}{D3869B}{\relax}
\gbDefTok{Char}{B8BB26}{\relax}
\gbDefTok{String}{B8BB26}{\relax}
\gbDefTok{VerbatimString}{B8BB26}{\relax}
\gbDefTok{SpecialString}{B8BB26}{\relax}
\gbDefTok{SpecialChar}{FE8019}{\relax}
\gbDefTok{Import}{FE8019}{\relax}
\gbDefTok{Preprocessor}{FE8019}{\relax}
\gbDefTok{Comment}{928374}{\textit}
\gbDefTok{Documentation}{928374}{\textit}
\gbDefTok{Annotation}{928374}{\textit}
\gbDefTok{CommentVar}{928374}{\textit}
\gbDefTok{Information}{928374}{\textit}
\gbDefTok{RegionMarker}{928374}{\relax}
\gbDefTok{Other}{8EC07C}{\relax}
\gbDefTok{Function}{8EC07C}{\relax}
\gbDefTok{Operator}{8EC07C}{\relax}
\gbDefTok{BuiltIn}{8EC07C}{\relax}
\gbDefTok{Extension}{8EC07C}{\relax}
\gbDefTok{Warning}{FE8019}{\textbf}
\gbDefTok{Alert}{FB4934}{\textbf}
\gbDefTok{Error}{FB4934}{\relax}

\newcommand{\gbnl}[1]{\llap{\textcolor{gbcodnumero}{#1}\hspace{1.2em}}}
\newcommand{\gbnc}{\relax}
\newcommand{\gbret}{\textcolor{gbcodnumero}{\symbol{"21A9}}}

\tcbset{gbcodigo/.style={breakable, colback=gbcodfondo,
  colframe=gbcodmarco, coltext=gbcodtexto, boxrule=0.4pt, arc=3pt,
  colbacktitle=gbcodcabeza, coltitle=gbcodrotulo, titlerule=0pt,
  fonttitle=\ttfamily\scriptsize, toptitle=1pt, bottomtitle=1pt,
  boxsep=0pt, left=3em, right=6pt, top=5pt, bottom=5pt,
  lefttitle=6pt, before skip=\baselineskip, after skip=\baselineskip}}
% \VerbatimEnvironment: fancyvrb BUSCA \end{gbCodigo} Y NO
% \end{Verbatim}. ES LA FORMA DOCUMENTADA DE ENVOLVER UN Verbatim.
\newenvironment{gbCodigo}[1]
  {\VerbatimEnvironment
   \begin{tcolorbox}[gbcodigo,title={#1}]%
   \begin{Verbatim}[commandchars=\\\{\},fontsize=\footnotesize]}
  {\end{Verbatim}\end{tcolorbox}}

% EL CONTADOR DE LOS LISTADOS (SC-42). ANTES SE DECLARABA EN EL
% CONTRATO DE LIBROS CON \newcounter{gbcodigo}[chapter]; LA CLASE
% article NO TIENE chapter. EL REINICIO POR CAPÍTULO LO AGREGA EL
% CONTRATO DE LIBROS CON \counterwithin*, QUE HACE LO MISMO QUE EL
% ARGUMENTO OPCIONAL (\@addtoreset) Y NO TOCA \thegbcodigo
\newcounter{gbcodigo}

% EL PIE VA PEGADO AL BLOQUE: SE DESCUENTA PARTE DEL after skip
\newcommand{\gbPieCodigo}[2]{%
  \par\nopagebreak\vspace{-.5\baselineskip}%
  \refstepcounter{gbcodigo}\label{#1}%
  {\small\noindent\textbf{Código~\thegbcodigo.}\space #2\par}%
  \vspace{.5\baselineskip}}


% ============================================================
% HOJAS DE CORTESÍA
% ------------------------------------------------------------
% Dos páginas en blanco al inicio —una hoja— que todo libro
% lleva. Entran en la paginación: son i y ii, y el sumario
% empieza en iii.
%
% \hbox{} ES IMPRESCINDIBLE: UNA PÁGINA SIN NINGÚN MATERIAL NO
% SE EMITE, ASÍ QUE UN \newpage SOBRE UNA PÁGINA VACÍA NO
% PRODUCE NADA Y EL CONTADOR NO AVANZA. LA CAJA VACÍA DA EL
% MÍNIMO DE CONTENIDO QUE OBLIGA A COMPONERLA.
%
% \thispagestyle{empty} LES QUITA EL FOLIO: SE CUENTAN PERO NO
% SE IMPRIME EL NÚMERO.
% ============================================================
\newcommand{\gbHojasCortesia}{%
	\thispagestyle{empty}\hbox{}\newpage
	\thispagestyle{empty}\hbox{}\newpage}

% ============================================================
% ENCABEZADO DE PIEZA GENERADA
% ------------------------------------------------------------
% Para las piezas sin cuerpo que llevan título propio —los tres
% índices, la bibliografía declarada como pieza—. El título sale
% de capitulos.titulo_capitulo, que es donde el editor lo
% escribe, y no del preámbulo.
%
% EL TÍTULO VA AL FOLIO CORRIDO POR \gbMarcasPieza, EL MISMO
% CAMINO QUE USAN LAS PIEZAS CON CUERPO. NO SE USA \markboth
% DIRECTO: ESO PONÍA EL TÍTULO EN LOS DOS LADOS DEL FOLIO Y
% CONVIVÍA MAL CON EL ESTILO DE PÁGINA, QUE YA LEE
% \gbTituloFolio Y \gbAutoriaFolio.
%
% LA AUTORÍA VA VACÍA: UN ÍNDICE NO TIENE AUTOR.
%
% \gbEncabezadoPieza{PLANO}{COMPUESTO} (CONTRATO 6): EL TÍTULO
% COMPUESTO VA A LA APERTURA; EL PLANO, AL SUMARIO Y AL FOLIO
% (SC-28).
% ============================================================
\newcommand{\gbEncabezadoPieza}[2]{%
	\chapter*{#2}%
	\addcontentsline{toc}{chapter}{#1}%
	\gbMarcasPieza{#1}{}}

% ============================================================
% ÍNDICE CON SU ENCABEZADO
% ------------------------------------------------------------
% \printindex HACE UN \clearpage INTERNO AUNQUE EL ÍNDICE SE
% DECLARE CON title={}. SIN NEUTRALIZARLO, EL ENCABEZADO QUEDA
% SOLO EN UNA PÁGINA Y LA LISTA EMPIEZA EN LA SIGUIENTE.
%
% EL \let VA DENTRO DE UN GRUPO: FUERA DE ÉL \clearpage TIENE QUE
% SEGUIR FUNCIONANDO, Y EL COLOFÓN QUE VIENE DESPUÉS LO NECESITA.
%
% \gbIndice{PLANO}{COMPUESTO}{ÍNDICE} (CONTRATO 6).
% ============================================================
\newcommand{\gbIndice}[3]{%
	\gbEncabezadoPieza{#1}{#2}%
	\begingroup
	\let\clearpage\relax
	\raggedright
	\small
	\printindex[#3]%
	\endgroup}

% ============================================================
% MARCAS DE PIEZA PARA EL FOLIO CORRIENTE
% ------------------------------------------------------------
% El folio lleva el título abreviado y la autoría de la pieza,
% que son datos de la base y no del manuscrito. El generador los
% emite antes de cada \include y el estilo de página los lee.
%
% \chaptertitle de titlesec trae el título COMPLETO, que en un
% folio de una línea no entra. Y la autoría del capítulo no
% existe en LaTeX: en el preámbulo legacy se declaraba a mano
% con \Author{} en cada pieza.
%
% SE DEFINEN VACÍAS PARA QUE EL ESTILO DE PÁGINA PUEDA LEERLAS
% ANTES DE QUE EL GENERADOR LAS COMPLETE: EN PRELIMINARES Y
% POSLIMINARES HAY PIEZAS QUE NO LLEVAN MARCAS.
% ============================================================
\newcommand{\gbTituloFolio}{}
\newcommand{\gbAutoriaFolio}{}
\newcommand{\gbMarcasPieza}[2]{%
  \renewcommand{\gbTituloFolio}{#1}%
  \renewcommand{\gbAutoriaFolio}{#2}}

% ============================================================
% AUTORÍA EN LA APERTURA DE LA PIEZA (CONTRATO 4)
% ------------------------------------------------------------
% En un libro colectivo, debajo del título de cada pieza van sus
% autores; en un libro de autor, solo en la pieza que firma otra
% persona, como un prólogo ajeno. El generador emite
% \gbAperturaPieza{LÍNEA} antes de cada \include: la línea ya
% viene armada —cada nombre en \gbNombreApertura, comas y «y» o
% «e» en letra normal— o vacía.
%
% \gbComponerApertura LA COMPONE CON \gbDisenoApertura Y LA
% VACÍA. LA ESTÉTICA LO LLAMA EN EL after-code DE
% \titleformat{\chapter}, QUE CORRE TAMBIÉN CON \chapter*.
% SE VACÍA AL USARLA PORQUE LA BIBLIOGRAFÍA DE biblatex Y LOS
% ÍNDICES TAMBIÉN ABREN CON \chapter* SIN PASAR POR EL
% GENERADOR: SIN ESO, REPETIRÍAN LOS AUTORES DE LA PIEZA
% ANTERIOR. \global PORQUE EL after-code CORRE DENTRO DE UN
% GRUPO.
% ============================================================
\def\gbLineaApertura{}
\newcommand{\gbAperturaPieza}[1]{\gdef\gbLineaApertura{#1}}
\newcommand{\gbComponerApertura}{%
  \ifx\gbLineaApertura\empty\else
    \gbDisenoApertura{\gbLineaApertura}%
    \global\let\gbLineaApertura\empty
  \fi}

% ============================================================
% AUTORÍA EN EL SUMARIO (CONTRATO 5)
% ------------------------------------------------------------
% En un libro colectivo —o en una pieza que firma otra persona
% en un libro de autor, como un prólogo ajeno— la entrada de la
% pieza en el sumario lleva arriba sus autores. El generador emite
% \gbIndicePieza{LÍNEA} antes de cada \include, con los nombres
% completos ya unidos, o vacía.
%
% EL DATO TIENE QUE QUEDAR EN EL .toc ANTES DE LA ENTRADA DE LA
% PIEZA. POR ESO SE ENVUELVE \addcontentsline: CUANDO LA ENTRADA
% ES UN CAPÍTULO DEL SUMARIO (toc/chapter), PRIMERO ESCRIBE
% \gbAutoresIndice{LÍNEA} Y DESPUÉS LA ENTRADA. PASAN POR ACÁ
% LOS DOS CASOS: \chapter, QUE ESCRIBE SU ENTRADA SOLO, Y
% \chapter* SEGUIDO DE \addcontentsline.
%
% LA LÍNEA SE VACÍA AL ESCRIBIRLA, COMO LA DE LA APERTURA: LA
% BIBLIOGRAFÍA Y LOS ÍNDICES NO SE QUEDAN CON AUTORES AJENOS.
%
% SE ENVUELVE EN \AtBeginDocument PARA TOMAR LA VERSIÓN DE
% \addcontentsline QUE DEJA hyperref, QUE SE CARGA DESPUÉS. UN
% SOLO # DENTRO DE \AtBeginDocument (VERIFICADO).
%
% AL LEER EL .toc, \gbAutoresIndice GUARDA LA LÍNEA Y LA ENTRADA
% SIGUIENTE LA COMPONE CON \gbComponerAutoresIndice, QUE LA
% ESTÉTICA LLAMA DESDE \titlecontents{chapter}.
% ============================================================
\def\gbLineaIndice{}
\def\gbIndiceTipoCapitulo{toc/chapter}
\newcommand{\gbIndicePieza}[1]{\gdef\gbLineaIndice{#1}}
\AtBeginDocument{%
  \let\gbAddContentsLineOriginal\addcontentsline
  \renewcommand{\addcontentsline}[3]{%
    \def\gbIndiceTipo{#1/#2}%
    \ifx\gbIndiceTipo\gbIndiceTipoCapitulo
      \ifx\gbLineaIndice\empty\else
        \addtocontents{toc}{\protect\gbAutoresIndice{\gbLineaIndice}}%
        \global\let\gbLineaIndice\empty
      \fi
    \fi
    \gbAddContentsLineOriginal{#1}{#2}{#3}}}
\def\gbAutoresEnIndice{}
\newcommand{\gbAutoresIndice}[1]{\gdef\gbAutoresEnIndice{#1}}
\newcommand{\gbComponerAutoresIndice}{%
  \ifx\gbAutoresEnIndice\empty\else
    \gbDisenoAutoresIndice{\gbAutoresEnIndice}%
    \global\let\gbAutoresEnIndice\empty
  \fi}

% ============================================================
% 1. LO DE LIBROS QUE PIDE babel: SIGLAS Y GLOSARIO (1.18.0)
% ------------------------------------------------------------
% LOS PAQUETES, babel Y csquotes ESTÁN EN EL CONTRATO COMÚN
% ============================================================

% EL DESARROLLO DE UNA SIGLA EN OTRO IDIOMA (1.18.0): CURSIVA Y LA
% PARTICIÓN DE SU IDIOMA. LO ESCRIBE docbook-to-latex.xsl DESDE UN
% <foreignphrase xml:lang="en">: \gbextranjero{en}{World Health
% Organization}. EL NOMBRE DE babel DE CADA CÓDIGO LO DEFINE GAMBAS
% (m_GenerarPDFLibro.DefinicionesCita) CON \csname gbbabel@<código>,
% Y AGREGA ESE IDIOMA A \gbidiomas PARA QUE babel LO CARGUE. UN CÓDIGO
% SIN DEFINIR CORTA LA COMPILACIÓN CON UN MENSAJE: COMPONER CON OTRA
% PARTICIÓN SERÍA UN ERROR QUE NADIE VE (MEDIDO CON lualatex Y babel
% DE TeX Live 2023: \foreignlanguage ACEPTA EL \csname, Y DENTRO
% \languagename ES english).
\newcommand*{\gbextranjero}[2]{%
  \ifcsname gbbabel@#1\endcsname
    \foreignlanguage{\csname gbbabel@#1\endcsname}{\emph{#2}}%
  \else
    \PackageError{gbpublisher}{El idioma «#1» de una sigla no tiene
      nombre de babel}{Las definiciones del proyecto no lo traen: hay
      que generar el PDF desde gbpublisher}%
  \fi}

% LA LISTA DE SIGLAS DEL LIBRO (1.18.0): DOS COLUMNAS ALINEADAS, LA
% SIGLA Y SU DESARROLLO (DECISIÓN DE ALBERTO, 31/10). LA ESCRIBE
% docbook-to-latex.xsl ASÍ:
%   \gbSiglasInicio
%   \gbSiglaMedir{OMS} \gbSiglaMedir{SITRAC-SITRAM} …
%   \begin{gbsiglas}
%   \item[OMS] Organización Mundial de la Salud
%   \end{gbsiglas}
% LA PRIMERA COLUMNA MIDE LO QUE LA SIGLA MÁS ANCHA DE LA LISTA; LA
% SEGUNDA EMPIEZA \gbSiglasSep DESPUÉS, Y UN DESARROLLO LARGO SIGUE
% ALINEADO EN ELLA. ES UNA list DE LaTeX: SE PARTE ENTRE PÁGINAS.
% \gbSiglasSep ES DE LA ESTÉTICA, QUE PUEDE CAMBIARLA CON \setlength.
\newlength{\gbSiglasSep}
\setlength{\gbSiglasSep}{1.5em}
\newlength{\gbSiglasAncho}
\newlength{\gbSiglaTmp}
\newcommand{\gbSiglasInicio}{\setlength{\gbSiglasAncho}{0pt}}
\newcommand{\gbSiglaMedir}[1]{%
  \settowidth{\gbSiglaTmp}{#1}%
  \ifdim\gbSiglaTmp>\gbSiglasAncho
    \setlength{\gbSiglasAncho}{\gbSiglaTmp}%
  \fi}
\newenvironment{gbsiglas}
  {\list{}{%
     \setlength{\labelwidth}{\gbSiglasAncho}%
     \setlength{\labelsep}{\gbSiglasSep}%
     \setlength{\leftmargin}{\gbSiglasAncho}%
     \addtolength{\leftmargin}{\gbSiglasSep}%
     \setlength{\itemsep}{0pt}%
     \setlength{\parsep}{0pt}%
     \setlength{\topsep}{0pt}%
     \renewcommand{\makelabel}[1]{##1\hfil}}}
  {\endlist}

% EL GLOSARIO DEL LIBRO (1.18.0): PÁRRAFO FRANCÉS (DECISIÓN DE
% ALBERTO, 31/10). EL TÉRMINO Y LA DEFINICIÓN EN EL MISMO PÁRRAFO; LAS
% LÍNEAS SIGUIENTES Y LOS DEMÁS PÁRRAFOS, SANGRADOS \gbGlosarioSangria.
% LO ESCRIBE docbook-to-latex.xsl ASÍ, UNA LETRA POR VEZ:
%   \gbGlosarioLetra{T}
%   \begin{gbglosario}
%   \item[]\gbGlosarioAncla{cap-9.gl-traveling}\gbGlosarioTermino{Traveling}%
%     \gbGlosarioVease{cap-9.gl-travelling}{Travelling}
%   \item[]\gbGlosarioAncla{…}\gbGlosarioTermino{\emph{Travelling}}\gbGlosarioSep
%     Desplazamiento … \gbGlosarioRemite{cap-9.gl-plano-secuencia}{plano secuencia}.
%   \end{gbglosario}
% EL ASPECTO —LA LETRA, LA NEGRITA, LA SEPARACIÓN, «véase»— ES DE LA
% ESTÉTICA, QUE PUEDE REDEFINIR ESTAS MACROS; ESTA CAPA GARANTIZA QUE
% EXISTAN. «véase» VA EN MINÚSCULA Y DESPUÉS DE COMA, COMO EN LOS
% ÍNDICES (DECISIÓN DE ALBERTO, 30/10, SC-77).
\newlength{\gbGlosarioSangria}
\setlength{\gbGlosarioSangria}{1.5em}
\newcommand{\gbGlosarioLetra}[1]{%
  \par\bigskip\noindent{\large\bfseries #1}\par\nopagebreak\smallskip}
\newcommand{\gbGlosarioTermino}[1]{\textbf{#1}}
% EL SEPARADOR ENTRE EL TÉRMINO Y LA DEFINICIÓN: SEMIRRAYA ENTRE
% ESPACIOS, EL PRIMERO DURO PARA QUE LA RAYA NO EMPIECE UNA LÍNEA
% (DECISIÓN DE ALBERTO, 31/10: «ONU – La Organización…»). EL CARÁCTER
% VA LITERAL Y NO COMO -- : NO DEPENDE DE LAS LIGADURAS DE LA FUENTE
\newcommand{\gbGlosarioSep}{\nobreakspace–\space}
\newcommand{\gbGlosarioAncla}[1]{\phantomsection\label{#1}}
\newcommand{\gbGlosarioRemite}[2]{\hyperref[#1]{#2}}
\newcommand{\gbGlosarioVease}[2]{, véase \hyperref[#1]{#2}}
\newenvironment{gbglosario}
  {\list{}{%
     \setlength{\leftmargin}{\gbGlosarioSangria}%
     \setlength{\itemindent}{-\gbGlosarioSangria}%
     \setlength{\labelwidth}{0pt}%
     \setlength{\labelsep}{0pt}%
     \setlength{\listparindent}{0pt}%
     \setlength{\itemsep}{0.4\baselineskip}%
     \setlength{\parsep}{0pt}%
     \setlength{\topsep}{0pt}}}
  {\endlist}



% ============================================================
% 3. BIBLIOGRAFÍA
% ------------------------------------------------------------
% EL ARCHIVO biblatex-<estandar>-config.tex ES RESPONSABLE DE
% CARGAR biblatex CON EL ESTILO QUE CORRESPONDA. SON NUESTROS
% ARCHIVOS, VIAJAN EN EL PAQUETE Y SE COPIAN AL PROYECTO: NO
% SE DEPENDE DE NINGÚN ESTILO REGISTRADO EN TeX Live.
%
% LA RAMA bibtex (vancouver) NO ESTÁ IMPLEMENTADA. EL XSLT
% CORTA ANTES DE GENERAR SI EL ESTILO LA REQUIERE.
%
% INMEDIATAMENTE DESPUÉS, LAS CADENAS EN CASTELLANO COMUNES A TODOS
% LOS ESTILOS (cita-cadenas-es.tex): NECESITAN biblatex YA CARGADO.
% ============================================================
\input{cita-\estandar-config.tex}
% EL REINICIO DE LAS CITAS POR CAPÍTULO (ibid., op. cit. Y LA PRIMERA
% CITA COMPLETA). ESTABA EN CADA cita-*-config.tex; SALIÓ DE AHÍ PORQUE
% ESOS ARCHIVOS LOS USAN TAMBIÉN LAS REVISTAS, CUYA CLASE NO TIENE
% chapter. DESPUÉS DE CARGAR biblatex LA OPCIÓN SE ACEPTA Y RIGE IGUAL
% (MEDIDO CON biblatex DE TeX Live 2023: EN EL SEGUNDO CAPÍTULO LA
% CITA VUELVE A SALIR COMPLETA; SIN ELLA SALE «ibid.»)
\ExecuteBibliographyOptions{citereset=chapter}
% ============================================================
% cita-cadenas-es.tex — CADENAS EN CASTELLANO COMUNES A TODOS LOS
% ESTILOS DE biblatex
% ------------------------------------------------------------
% SOLO LAS CADENAS QUE NI biblatex (spanish.lbx) NI EL ESTILO TRAEN:
% SON LA TRADUCCIÓN DE LOS CÓDIGOS DEL CATÁLOGO DE LA APLICACIÓN
% (TABLA cmb_biblatex: editor_type, thesis_type, pubstate) Y DE LOS
% AGREGADOS PROPIOS DE LA EDITORIAL. UNA CADENA ASÍ SOLO APARECE SI
% UNA ENTRADA LA USA: EN UN ESTILO QUE NO LA NECESITA NO CAMBIA NADA.
%
% LAS QUE spanish.lbx YA TRAE Y UN ESTILO REDEFINE (meses, url,
% phdthesis, compiler…) NO VAN ACÁ: SON DECISIONES DE CADA ESTILO Y
% QUEDAN EN SU cita-<estilo>-config.tex.
%
% LO CARGA QUIEN CARGA EL ESTILO, INMEDIATAMENTE DESPUÉS DE
% cita-<estilo>-config.tex (\NewBibliographyString EXIGE biblatex YA
% CARGADO):
%   libros  — preambulo-contrato.tex (PDF) Y EL DOCUMENTO MÍNIMO DE
%             preparar_citas_libro.sh (PASADA DE tex4ht), POR NOMBRE
%   revistas — m_XML, CON RUTA ABSOLUTA A ~/.gbpublisher/latex/
% NO LO CARGA EL CONFIG: \input RESUELVE CONTRA LA CARPETA DE LA
% COMPILACIÓN, NO CONTRA LA DEL ARCHIVO QUE LO PIDE, Y EN REVISTAS EL
% CONFIG SE CARGA POR RUTA ABSOLUTA DESDE OTRA CARPETA.
%
% \NewBibliographyString CON UNA CADENA QUE YA EXISTE NO HACE NADA
% (biblatex.sty); CON UNA NUEVA, LA DECLARA PARA TODOS LOS IDIOMAS.
% ============================================================

% SIN biblatex NO HAY CADENAS QUE DEFINIR: ES EL CASO DE vancouver,
% QUE VA POR bibtex (cita-vancouver-config.tex). EL ARCHIVO TERMINA
% ACÁ EN LUGAR DE CORTAR LA COMPILACIÓN CON \NewBibliographyString
% INDEFINIDO
\makeatletter
\@ifpackageloaded{biblatex}{\makeatother}{\makeatother\endinput}

% ------------------------------------------------------------
% 1. FUNCIONES DE EDITOR (editortype) QUE spanish.lbx NO TRAE
% ------------------------------------------------------------
% EL CATÁLOGO OFRECE coordinator Y direction. biblatex BUSCA LA
% CADENA DEL MISMO NOMBRE, SU PLURAL (…s) Y LA FORMA by…; SIN
% ELLAS IMPRIME EL CÓDIGO TAL CUAL («coordinator», MEDIDO EN IEEE).
% organizator ES UN ALIAS HEREDADO DE ISO 690: EL CATÁLOGO USA
% organizer, QUE spanish.lbx YA TRAE.
\NewBibliographyString{coordinator}
\NewBibliographyString{coordinators}
\NewBibliographyString{bycoordinator}
\NewBibliographyString{direction}
\NewBibliographyString{directions}
\NewBibliographyString{bydirection}
\NewBibliographyString{organizator}
\NewBibliographyString{organizators}
\NewBibliographyString{byorganizator}

% ------------------------------------------------------------
% 2. TIPOS DE TESIS Y TRABAJOS FINALES (thesis_type)
% ------------------------------------------------------------
% LOS CÓDIGOS DEL CATÁLOGO, SALVO mathesis, phdthesis Y candthesis,
% QUE spanish.lbx TRAE Y CADA ESTILO REDEFINE EN SU CONFIG. LOS
% ÚLTIMOS SEIS SON ALIAS HEREDADOS DE ISO 690, QUE EL CATÁLOGO NO
% OFRECE: QUEDAN PARA UN .bib ESCRITO A MANO.
\NewBibliographyString{bachelorthesis}
\NewBibliographyString{engineerthesis}
\NewBibliographyString{specializationthesis}
\NewBibliographyString{fieldwork}
\NewBibliographyString{finalessay}
\NewBibliographyString{finalmonograph}
\NewBibliographyString{finalproject}
\NewBibliographyString{gradothesis}
\NewBibliographyString{integrativefinalwork}
\NewBibliographyString{magisterthesis}
\NewBibliographyString{masterfinalwork}
\NewBibliographyString{postgradthesis}
\NewBibliographyString{teacherthesis}
\NewBibliographyString{researchproject}
\NewBibliographyString{residencyreport}
\NewBibliographyString{shortthesis}
\NewBibliographyString{technicalreport}
\NewBibliographyString{undergraduatethesis}
\NewBibliographyString{ingenieriathesis}
\NewBibliographyString{licenciaturathesis}
\NewBibliographyString{maestriathesis}
\NewBibliographyString{posgradothesis}
\NewBibliographyString{profesoradothesis}
\NewBibliographyString{tifthesis}

% ------------------------------------------------------------
% 3. ESTADOS DE PUBLICACIÓN (pubstate) Y FORMAS DE EDICIÓN
% ------------------------------------------------------------
\NewBibliographyString{mimeograph}
\NewBibliographyString{photocopy}
\NewBibliographyString{digitalprinting}
\NewBibliographyString{edauthor}
% self-published ES EL CÓDIGO DEL CATÁLOGO PARA LA EDICIÓN DE AUTOR (SIN
% CADENA, biblatex LO IMPRIMÍA TAL CUAL: MEDIDO EN VERONA E IEEE). edauthor
% QUEDA PARA UN .bib ESCRITO A MANO. EL GUION EN LA CLAVE ES VÁLIDO (MEDIDO)
\NewBibliographyString{self-published}
\NewBibliographyString{documentjob}
\NewBibliographyString{origpubbare}

% ------------------------------------------------------------
% 4. CINE: DIRECCIÓN, GUION Y ELENCO
% ------------------------------------------------------------
\NewBibliographyString{dirigida}
\NewBibliographyString{dirigidas}
\NewBibliographyString{bydirigida}
\NewBibliographyString{escrita}
\NewBibliographyString{escritas}
\NewBibliographyString{byescrita}
\NewBibliographyString{elenco}
\NewBibliographyString{elencos}
\NewBibliographyString{byelenco}

% ------------------------------------------------------------
% 5. «Y SIGUIENTE», «Y SIGUIENTES»: \sq Y \sqq EN UNA postnote
% ------------------------------------------------------------
\NewBibliographyString{sq}
\NewBibliographyString{sqq}
\providecommand*{\sq}{\bibstring{sq}}
\providecommand*{\sqq}{\bibstring{sqq}}

% ------------------------------------------------------------
% 6. EL TEXTO EN CASTELLANO
% ------------------------------------------------------------
\DefineBibliographyStrings{spanish}{%
  coordinator          = {coord\adddot},
  coordinators         = {coords\adddot},
  bycoordinator        = {coord\adddot\addspace por},
  direction            = {dir\adddot},
  directions           = {dirs\adddot},
  bydirection          = {dir\adddot\addspace por},
  organizator          = {org\adddot},
  organizators         = {orgs\adddot},
  byorganizator        = {org\adddot\addspace por},
  bachelorthesis       = {Tesis de Licenciatura},
  engineerthesis       = {Tesis de Ingeniería},
  specializationthesis = {Tesis de Especialización},
  fieldwork            = {Trabajo de Campo},
  finalessay           = {Ensayo Final},
  finalmonograph       = {Monografía Final de Carrera},
  finalproject         = {Proyecto Final},
  gradothesis          = {Tesis de Grado},
  integrativefinalwork = {Trabajo Integrador Final (TIF)},
  magisterthesis       = {Tesis de Magíster},
  masterfinalwork      = {Trabajo Final de Maestría},
  postgradthesis       = {Tesis de Posgrado},
  teacherthesis        = {Tesis de Profesorado},
  researchproject      = {Proyecto de Investigación},
  residencyreport      = {Informe Final de Residencia},
  shortthesis          = {Tesina},
  technicalreport      = {Trabajo Final Técnico},
  undergraduatethesis  = {Trabajo Final de Grado},
  ingenieriathesis     = {Tesis de Ingeniería},
  licenciaturathesis   = {Tesis de Licenciatura},
  maestriathesis       = {Tesis de Maestría},
  posgradothesis       = {Tesis de Posgrado},
  profesoradothesis    = {Tesis de Profesorado},
  tifthesis            = {Trabajo Integrador Final},
  mimeograph           = {mimeo},
  photocopy            = {fotocopia},
  digitalprinting      = {impresión digital},
  edauthor             = {edición de autor},
  self-published       = {edición de autor},
  documentjob          = {documento de trabajo},
  origpubbare          = {orig\adddotspace pub\adddotspace},
  dirigida             = {dirección},
  dirigidas            = {dirección},
  bydirigida           = {Dirección de},
  escrita              = {escrita},
  escritas             = {escrita},
  byescrita            = {escrita por},
  elenco               = {actúan:},
  elencos              = {actúan:},
  byelenco             = {actúan:},
  sq                   = {~y~sig\adddot},
  sqq                  = {~y~ss\adddot},
}


% ============================================================
% Índice de autores del aparato bibliográfico
% ------------------------------------------------------------
% Lo llena biblatex: con indexing=true en el estilo, cada nombre
% citado y cada nombre de la bibliografía se indexan solos. El
% editor no escribe ningún \index.
%
% esindex ORDENA LOS APELLIDOS COMPUESTOS CASTELLANOS COMO
% CORRESPONDE: «de la Torre» SE ALFABETIZA POR Torre PERO SE
% IMPRIME CON LA PARTÍCULA. SIN ÉL, EL ORDEN POR OMISIÓN LOS
% MANDA A LA D.
%
% VA EN EL CONTRATO PORQUE ES LO QUE HACE QUE LA PIEZA
% indice_autores_bibliografia DIGA LA VERDAD. EL \makeindex ES
% DISEÑO Y VA EN LA ESTÉTICA.
% ============================================================
% LOS CARACTERES ! @ | SON SINTAXIS DEL ÍNDICE (NIVEL, CLAVE Y FORMATO):
% UN NOMBRE QUE LOS CONTIENE, COMO «Democracy Now!», DEJA UNA CLAVE
% VACÍA Y xindy ABORTA EL ÍNDICE ENTERO («CHAR: index 0 should be less
% than the length of the string»). SE ESCAPAN CON LA COMILLA DE xindy,
% ENTRE LLAVES PARA QUE esindex NO PARTA EN ELLOS. SOLO AFECTA AL
% ÍNDICE: LA BIBLIOGRAFÍA IMPRIME EL NOMBRE TAL CUAL. ESCAPARLO EN EL
% .bib NO SIRVE: CON babel spanish, "! SE IMPRIME COMO ¡.
\ExplSyntaxOn
\cs_new_protected:Npn \gb_escapar_indice:N #1
  {
    \tl_replace_all:Nnn #1 { ! } { {"!} }
    \tl_replace_all:Nnn #1 { @ } { {"@} }
    \tl_replace_all:Nnn #1 { | } { {"|} }
  }
\cs_new_protected:Npn \gbEscaparIndice
  {
    \gb_escapar_indice:N \namepartfamily
    \gb_escapar_indice:N \namepartgiven
    \gb_escapar_indice:N \namepartprefix
    \gb_escapar_indice:N \namepartsuffix
  }
\ExplSyntaxOff
\usepackage{esindex}
% EL FORMATO DE PÁGINA VA EXPLÍCITO: CON hyperindex=false (GV-159)
% hyperref YA NO LO AGREGA, Y SIN \hyperpage LA PÁGINA SALE SIN
% ENLACE. esindex ACEPTA |formato (MEDIDO: LOS ENLACES SE CONSERVAN).
\newcommand{\gbIndiceNombre}[1]{\esindex[names]{#1|hyperpage}}
\DeclareIndexNameFormat{default}{%
	\gbEscaparIndice
	\usebibmacro{index:name}{\gbIndiceNombre}
	{\namepartfamily}
	{\namepartgiven}
	{\namepartprefix}
	{\namepartsuffix}}

% ============================================================
% ÍNDICES ONOMÁSTICO Y DE CONCEPTOS (CONTRATO 18, SC-77)
% ------------------------------------------------------------
% LOS CARGA EL EDITOR CON \index[onomastico]{…} Y
% \index[concepto]{…} EN EL .md; docbook-to-latex.xsl LOS
% VUELVE A ESCRIBIR CON EL FORMATO DE PÁGINA EXPLÍCITO:
%   |hyperpage        LA PÁGINA, CON ENLACE (hyperref)
%   |gbIndNota        LA MARCA ESTÁ EN UNA NOTA: «23 n.»
%   |gbIndPrincipal   LA PÁGINA PRINCIPAL, EN NEGRITA
% CADA FORMATO TIENE QUE ESTAR DECLARADO EN gbindice.xdy: SI NO,
% xindy DESCARTA LA PÁGINA CON UN WARNING Y NADA MÁS (GV-160).
% EL ESPACIO ANTES DE «n.» ES IRROMPIBLE: LA PÁGINA Y SU MARCA NO
% SE SEPARAN EN DOS LÍNEAS.
% ============================================================
\newcommand{\gbIndNota}[1]{\hyperpage{#1}~n.}
\newcommand{\gbIndPrincipal}[1]{\textbf{\hyperpage{#1}}}

% ============================================================
% 4. PÁGINAS Y SALTOS
% ============================================================

% ------------------------------------------------------------
% \gbclearrecto — SALTA A LA PRÓXIMA PÁGINA IMPAR
% \gbclearverso — SALTA A LA PRÓXIMA PÁGINA PAR
% ------------------------------------------------------------
% LA PÁGINA DE RELLENO LLEVA \hbox{} PORQUE UNA PÁGINA SIN
% CONTENIDO NO SE EMITE. NO SE USA UN CARÁCTER BLANCO: ESO
% DEJA TEXTO EXTRAÍBLE EN EL PDF, QUE APARECE AL COPIAR, EN
% pdftotext Y EN UN LECTOR DE PANTALLA.
% EL CONTADOR SIGUE CORRIENDO: SOLO SE OCULTA EL FOLIO.
% ------------------------------------------------------------
\makeatletter
\newcommand{\gbclearrecto}{%
  \clearpage
  \ifodd\c@page\else
    \thispagestyle{empty}\hbox{}\newpage
  \fi}
\newcommand{\gbclearverso}{%
  \clearpage
  \ifodd\c@page
    \thispagestyle{empty}\hbox{}\newpage
  \fi}
% LA BLANCA QUE GENERA \cleardoublepage SALE, POR OMISIÓN, CON
% EL FOLIO Y EL ENCABEZADO CORRIENTES. ESTE PARCHE LA DEJA
% VACÍA, QUE ES LO QUE SE ESPERA DE UNA PÁGINA EN BLANCO.
\renewcommand{\cleardoublepage}{\gbclearrecto}
\makeatother

% ------------------------------------------------------------
% gbpaginaespecial / gbpaginaimagen
% ------------------------------------------------------------
% PIEZAS SUELTAS DE PRELIMINARES: UNA DEDICATORIA, UN
% RECORDATORIO, UNA FOTO A PÁGINA ENTERA. VAN EN RECTO, SIN
% FOLIO NI ENCABEZADO, Y SU VUELTA QUEDA EN BLANCO. NO LLEVAN
% \chapter Y NO ENTRAN AL ÍNDICE.
% LA DE IMAGEN ROMPE LA CAJA DE TEXTO; LA OTRA LA RESPETA.
% ------------------------------------------------------------
\newenvironment{gbpaginaespecial}
  {\gbclearrecto\thispagestyle{empty}\begin{center}}
  {\end{center}\gbclearrecto}

\newenvironment{gbpaginaimagen}
  {\gbclearrecto\thispagestyle{empty}%
   \begingroup\centering\setlength{\parindent}{0pt}}
  {\par\endgroup\gbclearrecto}

% ------------------------------------------------------------
% gbdedicatoria
% ------------------------------------------------------------
% LA DEDICATORIA ES UNA PIEZA SUELTA (fm-00 A fm-09) CON
% DISEÑO PROPIO: LA PÁGINA DE gbpaginaespecial —RECTO, SIN
% FOLIO NI ENCABEZADO, VUELTA EN BLANCO, SIN TÍTULO NI
% ENTRADA EN EL SUMARIO— Y ADENTRO EL BLOQUE DEL EPÍGRAFE SIN
% FILETE, PORQUE NO HAY SEGUNDA PARTE (\gbepigrafe CON #2
% VACÍO): 60 % DE LA COLUMNA, A LA DERECHA, \small, TEXTO EN
% BANDERA Y SIN CORTE DE PALABRA.
% LA BAJADA DESDE EL TOPE DE CAJA ES \gbBajadaDedicatoria, DE
% LA ESTÉTICA. EL \hbox{} ANTES DEL \vspace: UN ESPACIO
% VERTICAL AL PRINCIPIO DE UNA PÁGINA SE DESCARTA SI NO HAY
% NADA ANTES (EL MISMO RECURSO QUE \gbDisenoPortadilla).
% LA MINIPAGE VA CON [t]: SU PRIMERA LÍNEA QUEDA EN LA LÍNEA
% DEL PÁRRAFO QUE LA CONTIENE, Y LA BAJADA SE MIDE IGUAL CON
% UNA LÍNEA DE DEDICATORIA QUE CON CINCO.
% GEMELA DEL BLOQUE DE \gbepigrafe: UN CAMBIO DE DISEÑO DEL
% EPÍGRAFE VA TAMBIÉN ACÁ.
% ------------------------------------------------------------
\newenvironment{gbdedicatoria}
  {\gbclearrecto\thispagestyle{empty}%
   \hbox{}\vspace{\gbBajadaDedicatoria}%
   \par\noindent\hfill
   \begin{minipage}[t]{.6\linewidth}%
     \small\parindent=0pt\parskip=.5em
     \hyphenpenalty=10000\exhyphenpenalty=10000
     \raggedright}
  {\par\end{minipage}\par\gbclearrecto}

% ------------------------------------------------------------
% gbcolofon
% ------------------------------------------------------------
% VA EN VERSO. NO PUEDE SER UN \chapter*: EN LA CLASE book,
% \chapter FUERZA EL RECTO PORQUE @openright ESTÁ ACTIVO, Y
% CUALQUIER SALTO PREVIO QUEDA ANULADO.
% ------------------------------------------------------------
\newenvironment{gbcolofon}
  {\gbclearverso\thispagestyle{empty}\justifying}
  {\clearpage}

% ------------------------------------------------------------
% LAS PRIMERAS DEL LIBRO (SC-25)
% ------------------------------------------------------------
% PIEZAS SIN CUERPO: EL GENERADOR EMITE LA MACRO Y LOS DATOS
% ESTÁN EN preambulo-proyecto.tex. EL CONTRATO FIJA LA PÁGINA;
% EL ASPECTO ESTÁ EN LA ESTÉTICA, EN \gbDisenoPortadilla,
% \gbDisenoPortada Y \gbDisenoCreditos.
%
%   portadilla  IMPAR (iii, DESPUÉS DE LAS HOJAS DE CORTESÍA)
%   portada     IMPAR (v); SU VUELTA EN BLANCO LA PONE
%               \gbclearrecto
%   creditos    PAR, LA VUELTA DE LA PORTADA (vi)
%
% SIN FOLIO NI ENCABEZADO, Y FUERA DEL SUMARIO: NO USAN
% \chapter, QUE ADEMÁS FORZARÍA EL RECTO (VER gbcolofon).
% NO SE USA gbpaginaespecial: DEJA LA VUELTA EN BLANCO, Y LA
% VUELTA DE LA PORTADA SON LOS CRÉDITOS.
% ------------------------------------------------------------
\newcommand{\gbPortadilla}{%
  \gbclearrecto\thispagestyle{empty}%
  \gbDisenoPortadilla
  \clearpage}
\newcommand{\gbPortada}{%
  \gbclearrecto\thispagestyle{empty}%
  \gbDisenoPortada
  \clearpage}
\newcommand{\gbCreditos}{%
  \gbclearverso\thispagestyle{empty}%
  \gbDisenoCreditos
  \clearpage}

% SI UNA ESTÉTICA DE EDITORIAL NO DEFINE EL DISEÑO, EL LIBRO
% COMPILA IGUAL CON UNO MÍNIMO. SE RESUELVE AL EMPEZAR EL
% DOCUMENTO, CUANDO LA ESTÉTICA YA SE CARGÓ: \providecommand NO
% TOCA LO QUE ELLA HAYA DEFINIDO.
\AtBeginDocument{%
  % UN SOLO #: \AtBeginDocument NO ES UNA DEFINICIÓN, Y ## DA
  % «macro parameter character» (VERIFICADO)
  \providecommand{\gbNombrePortada}[1]{#1}%
  \providecommand{\gbNombreApertura}[1]{#1}%
  \providecommand{\gbDisenoAutoresIndice}[1]{#1\newline}%
  \providecommand{\gbDisenoApertura}[1]{%
    \begin{center}#1\end{center}}%
  % EL SUBTÍTULO CIERRA EL PÁRRAFO DEL TÍTULO Y ABRE EL SUYO, Y LO
  % CIERRA DENTRO DEL GRUPO DE \normalsize: EL INTERLINEADO DE UN
  % PÁRRAFO LO FIJA EL TAMAÑO VIGENTE AL CERRARLO. CON \\ EL
  % SUBTÍTULO QUEDABA EN EL PÁRRAFO DEL TÍTULO, Y SUS LÍNEAS SALÍAN
  % CON EL INTERLINEADO DE \LARGE. \endgraf Y NO \par: EL
  % ARGUMENTO DE \part NO ADMITE \par (VERIFICADO).
  \providecommand{\gbDisenoSubtituloCapitulo}[1]{%
    \endgraf\vspace{0.3em}{\normalsize\mdseries\itshape #1\endgraf}}%
  \providecommand{\gbDisenoPortadilla}{%
    \begin{center}\LARGE\gbtitulo\par\end{center}}%
  \providecommand{\gbDisenoPortada}{%
    \begin{center}\gbautoriaportada\par
    \ifx\gbleyendaautoria\empty\else\gbleyendaautoria\par\fi\bigskip
    {\LARGE\gbtitulo\par}\bigskip
    \gbsubtitulo\par\end{center}}%
  \providecommand{\gbBajadaDedicatoria}{30mm}%
  \providecommand{\gbDisenoCreditos}{%
    \noindent\gbTextoCatalogacion\par\vfill
    \noindent\gbTextoLegal\par}}

% PANDOC ESCRIBE \tightlist EN LAS LISTAS DE LOS TEXTOS DE
% CRÉDITOS. SU PLANTILLA LO DEFINE; UN FRAGMENTO NO LO TRAE.
\providecommand{\tightlist}{%
  \setlength{\itemsep}{0pt}\setlength{\parskip}{0pt}}

% ============================================================
% 5. ESTRUCTURA
% ============================================================

% ------------------------------------------------------------
% \gbContadoresSinCapitulo
% ------------------------------------------------------------
% SE LLAMA UNA VEZ AL ENTRAR EN PRELIMINARES Y UNA VEZ AL
% ENTRAR EN POSLIMINARES. NO POR PIEZA.
%
% EN ESAS DOS ZONAS \chapter NO EJECUTA \refstepcounter{chapter}
% (book.cls, \@chapter, RAMA \if@mainmatter). ESE COMANDO ES LO
% ÚNICO QUE INCREMENTA EL CONTADOR DE CAPÍTULO Y, A LA VEZ,
% REINICIA LOS QUE DEPENDEN DE ÉL: section, footnote, figure,
% table Y equation. SI NO CORRE, NO PASA NINGUNA DE LAS DOS.
%
% LLEVAR chapter A CERO ACTIVA LA GUARDA \ifnum \c@chapter>\z@
% DE \thefigure, \thetable Y \theequation, QUE OMITE EL PREFIJO.
% SIN ESTO, UNA FIGURA DE POSLIMINARES SALE "Figura 1.2", CON
% EL NÚMERO DEL ÚLTIMO CAPÍTULO DEL CUERPO. VERIFICADO.
%
% EL REINICIO POR PIEZA DE NOTAS, FIGURAS, CUADROS Y ECUACIONES
% NO VA ACÁ: LO HACE \gbCapituloSinNumero EN CADA PIEZA SIN
% NÚMERO (CONTRATO 9).
% ------------------------------------------------------------
\newcommand{\gbContadoresSinCapitulo}{%
  \setcounter{chapter}{0}%
  \setcounter{section}{0}%
  \setcounter{figure}{0}%
  \setcounter{table}{0}%
  \setcounter{equation}{0}%
  \setcounter{gbcodigo}{0}%
}

% ------------------------------------------------------------
% NUMERACIÓN POR PIEZA (CONTRATO 9)
% ------------------------------------------------------------
% REGLA: EN UN CAPÍTULO NUMERADO, LAS FIGURAS, LOS CUADROS Y LAS
% ECUACIONES VAN POR CAPÍTULO (2.1, 2.2; A.1 EN LOS APÉNDICES),
% COMO LOS NUMERA LA CLASE. EN TODA PIEZA SIN NÚMERO —DE
% PRELIMINARES, DEL CUERPO O DE POSLIMINARES— REINICIAN EN 1
% DENTRO DE LA PIEZA: 1, 2, 3. EL EPUB Y EL HTML REPITEN
% EXACTAMENTE ESTE NÚMERO (SC-31).
%
% \gbReiniciarPieza LO LLAMA \gbCapituloSinNumero, QUE ABRE
% TODA PIEZA SIN NÚMERO EN CUALQUIER ZONA: ASÍ LA REGLA VIVE EN
% UN SOLO LUGAR Y NO DEPENDE DE QUE EL main.tex SEPA EN QUÉ ZONA
% ESTÁ CADA TIPO. REINICIA TAMBIÉN LA NOTA AL PIE, QUE YA IBA
% POR PIEZA. UN CAPÍTULO NUMERADO NO LO NECESITA: \chapter
% REINICIA SUS CONTADORES.
%
% DOS PIEZAS CON SU «Figura 1» NO CHOCAN EN \label: LAS
% ETIQUETAS SON LOS xml:id (fig-mapa). LO QUE CHOCA ES EL ANCLA
% INTERNA DE hyperref, \theHfigure, QUE VALE \theHchapter.N:
% CON chapter EN CERO LAS DOS FIGURAS QUEDAN «figure.0.1», EL
% REGISTRO DICE «ignoring duplicate destination» Y EL ENLACE A
% LA SEGUNDA LLEVA A LA PRIMERA (VERIFICADO CON book Y hyperref
% SOLOS). LA ESTÉTICA ACTUAL CARGA caption, CUYAS ANCLAS
% (figure.caption.N) YA SON ÚNICAS: EL AJUSTE ESTÁ ACÁ PARA QUE
% LA GARANTÍA NO DEPENDA DE LA ESTÉTICA QUE PONGA UNA EDITORIAL.
% EL CONTADOR gbpieza, QUE AVANZA EN CADA REINICIO, ENTRA EN EL
% ANCLA Y LA HACE ÚNICA. SE REGISTRA EN EL GANCHO DE CARGA DE
% hyperref (COMO \gbCorteTitulo): hyperref SE CARGA EN LA
% ÚLTIMA CAPA.
%
% EN EL CUERPO EL CONTADOR DE CAPÍTULO NO SE TOCA (VER
% \gbPiezaSinNumerar), ASÍ QUE LA GUARDA DE LA CLASE NO
% ALCANZA: UNA FIGURA DE LAS CONCLUSIONES SALDRÍA «3.1». POR ESO
% \thefigure, \thetable Y \theequation CONSULTAN ADEMÁS
% \ifgbPiezaSinNumero, QUE ENCIENDE \gbCapituloSinNumero Y APAGA
% \gbCapitulo. EL RESTO DE LA DEFINICIÓN ES LA DE book.cls.
% VERIFICADO: PRÓLOGO 1, INTRODUCCIÓN 1, CAPÍTULO 1.1,
% CONCLUSIONES 1, APÉNDICE A.1, SIN ANCLAS REPETIDAS.
% ------------------------------------------------------------
\newcounter{gbpieza}
% EL CONTADOR DE LOS LISTADOS DE CÓDIGO (SC-42) LO DECLARA EL
% CONTRATO COMÚN; ACÁ SE LE SUMA EL REINICIO CON EL CAPÍTULO, COMO
% figure EN book.cls, Y SU \the SIGUE LA MISMA REGLA
\counterwithin*{gbcodigo}{chapter}
\newif\ifgbPiezaSinNumero
\newcommand{\gbReiniciarPieza}{%
  \stepcounter{gbpieza}%
  \setcounter{footnote}{0}%
  \setcounter{figure}{0}%
  \setcounter{table}{0}%
  \setcounter{equation}{0}%
  \setcounter{gbcodigo}{0}}
\makeatletter
\renewcommand{\thefigure}{%
  \ifgbPiezaSinNumero\else\ifnum\c@chapter>\z@\thechapter.\fi\fi
  \@arabic\c@figure}
\renewcommand{\thetable}{%
  \ifgbPiezaSinNumero\else\ifnum\c@chapter>\z@\thechapter.\fi\fi
  \@arabic\c@table}
\renewcommand{\theequation}{%
  \ifgbPiezaSinNumero\else\ifnum\c@chapter>\z@\thechapter.\fi\fi
  \@arabic\c@equation}
\renewcommand{\thegbcodigo}{%
  \ifgbPiezaSinNumero\else\ifnum\c@chapter>\z@\thechapter.\fi\fi
  \@arabic\c@gbcodigo}
\makeatother
\AddToHook{package/hyperref/after}{%
  \renewcommand{\theHfigure}{\theHchapter.\arabic{gbpieza}.\arabic{figure}}%
  \renewcommand{\theHtable}{\theHchapter.\arabic{gbpieza}.\arabic{table}}%
  \renewcommand{\theHequation}{\theHchapter.\arabic{gbpieza}.\arabic{equation}}%
  % hyperref NO DEFINE \theH DE UN CONTADOR PROPIO HASTA QUE SE USA:
  % \def Y NO \renewcommand, QUE FALLA CON «undefined» (MEDIDO)
  \def\theHgbcodigo{\theHchapter.\arabic{gbpieza}.\arabic{gbcodigo}}}

% ------------------------------------------------------------
% \gbPiezaSinNumerar
% ------------------------------------------------------------
% PARA UNA PIEZA SIN NUMERAR QUE VA EN LA ZONA DE CUERPO. EL
% CASO TÍPICO SON LAS CONCLUSIONES, QUE DEBEN PRECEDER A
% \appendix PARA QUE LOS APÉNDICES NUMEREN CON LETRA.
%
% COMO ESAS PIEZAS SE EMITEN CON \chapter*, NO CORRE
% \refstepcounter{chapter} Y NO SE REINICIA NADA: SU PRIMERA
% NOTA AL PIE ARRANCA DONDE LA DEJÓ EL CAPÍTULO ANTERIOR.
% VERIFICADO.
%
% NO SE TOCA EL CONTADOR DE CAPÍTULO: EL APÉNDICE POSTERIOR LO
% NECESITA INTACTO, Y \appendix YA LO LLEVA A CERO POR SU
% CUENTA. FIGURAS, CUADROS Y ECUACIONES LOS REINICIA
% \gbCapituloSinNumero (CONTRATO 9, ARRIBA).
% ------------------------------------------------------------
% TAMBIÉN APAGA LA NUMERACIÓN DE SECCIONES. SIN ESO, LA
% PRIMERA SECCIÓN DE LAS CONCLUSIONES SALE "1.1", QUE COLISIONA
% CON LA SECCIÓN 1.1 REAL DEL CAPÍTULO 1: \thesection ES
% \thechapter.\arabic{section} Y EL CONTADOR DE CAPÍTULO SIGUE
% EN 1. VERIFICADO EN EL SUMARIO.
% LA ZONA DE APÉNDICES LA VUELVE A ENCENDER.
\newcommand{\gbPiezaSinNumerar}{%
  \setcounter{footnote}{0}%
  \setcounter{section}{0}%
  \setcounter{secnumdepth}{0}%
}

% ------------------------------------------------------------
% \gbParte{título plano}{título compuesto}{subtítulo}
% ------------------------------------------------------------
% LAS PARTES NO SON ARCHIVOS: SON METADATOS DE partes_libro Y
% LAS EMITE GAMBAS EN EL main.tex. NO LLEVAN TEXTO PROPIO.
% EL SUBTÍTULO PUEDE VENIR VACÍO. AL ÍNDICE VA SOLO EL TÍTULO,
% EN SU VERSIÓN PLANA (CONTRATO 6, SC-28).
% EL SUBTÍTULO VA EN SU PROPIO PÁRRAFO, CERRADO DENTRO DE
% \normalsize, PARA QUE TOME SU INTERLINEADO Y NO EL DEL TÍTULO
% (VER \gbDisenoSubtituloCapitulo). \endgraf PORQUE EL ARGUMENTO
% DE \part NO ADMITE \par.
% ------------------------------------------------------------
\newcommand{\gbParte}[3]{%
  \part[#1]{#2%
    \ifstrempty{#3}{}{\endgraf\vspace{0.4em}{\normalsize\mdseries\itshape #3\endgraf}}%
  }}

% ------------------------------------------------------------
% \gbCapitulo{título plano}{título compuesto}{subtítulo}
% \gbCapituloSinNumero{título plano}{título compuesto}{subtítulo}
% ------------------------------------------------------------
% LAS EMITE docbook-to-latex.xsl PARA ABRIR CADA PIEZA CON
% CUERPO (CONTRATO 6). LA COMPUESTA LLEVA EL CORTE DEL EDITOR Y
% VA A LA APERTURA; LA PLANA VA AL SUMARIO Y A LOS MARCADORES.
% ES EL PATRÓN \chapter[sumario]{título}.
%
% LA SIN NÚMERO ES \chapter* MÁS \addcontentsline: ENTRA AL
% SUMARIO SIN NÚMERO. LO QUE ESCRIBE EN EL .toc PASA POR EL
% \addcontentsline ENVUELTO DE LA AUTORÍA (CONTRATO 5).
%
% EL SUBTÍTULO PUEDE VENIR VACÍO Y NO VA AL SUMARIO. SU ASPECTO
% LO DA \gbDisenoSubtituloCapitulo, QUE LA ESTÉTICA PUEDE
% REDEFINIR; SI NO LO HACE, VALE EL DE ABAJO.
% ------------------------------------------------------------
\newcommand{\gbTituloConSubtitulo}[2]{%
  #1\ifstrempty{#2}{}{\gbDisenoSubtituloCapitulo{#2}}}
% \gbCapitulo APAGA \ifgbPiezaSinNumero: UN CAPÍTULO NUMERADO
% DESPUÉS DE UNA PIEZA SIN NÚMERO DEL CUERPO VUELVE A LLEVAR EL
% PREFIJO (CONTRATO 9).
\newcommand{\gbCapitulo}[3]{%
  \global\gbPiezaSinNumerofalse
  \chapter[#1]{\gbTituloConSubtitulo{#2}{#3}}}
% \gbCapituloSinNumero REINICIA LOS CONTADORES DE LA PIEZA Y
% ENCIENDE \ifgbPiezaSinNumero (CONTRATO 9).
\newcommand{\gbCapituloSinNumero}[3]{%
  \gbReiniciarPieza
  \global\gbPiezaSinNumerotrue
  \chapter*{\gbTituloConSubtitulo{#2}{#3}}%
  \addcontentsline{toc}{chapter}{#1}}

% ============================================================
% 9. COLOFÓN
% ------------------------------------------------------------
% LAS FUENTES Y LA VERSIÓN SON VALORES, NO TEXTO ESCRITO. SI
% UNA EDITORIAL REEMPLAZA LA ESTÉTICA Y CAMBIA LAS FUENTES,
% EL COLOFÓN SIGUE DICIENDO LA VERDAD SIN QUE NADIE SE
% ACUERDE DE EDITARLO.
% LA FECHA ES LA DE PUBLICACIÓN Y NO LA DEL DÍA: UNA MARCA DE
% TIEMPO HARÍA QUE CADA COMPILACIÓN PRODUJERA UN DIFF AUNQUE
% NADA HUBIERA CAMBIADO.
% ============================================================
% SI LA ESTÉTICA NO LAS DECLARA, EL COLOFÓN LO DICE EN VEZ DE DAR UN
% "Undefined control sequence" QUE NO EXPLICA NADA.
\providecommand{\gbfuentetexto}{(sin declarar)}
\providecommand{\gbfuentetitulos}{(sin declarar)}
\providecommand{\gbversion}{(sin version)}

% EL RECUADRO SE DEFINE COMO ESTILO Y NO EN LÍNEA: ASÍ UNA EDITORIAL
% QUE QUIERA OTRO ASPECTO REDEFINE gbcolofoncaja SIN TOCAR EL TEXTO,
% QUE ES LO QUE EL CONTRATO GARANTIZA.
\tcbset{gbcolofoncaja/.style={
  colback=black!7,
  colframe=black!50,
  boxrule=0.7pt,
  arc=3mm,
  left=8pt, right=8pt, top=8pt, bottom=8pt,
  boxsep=0pt}}

\newcommand{\gbTextoColofon}{%
  \begin{tcolorbox}[gbcolofoncaja]
    La composición tipográfica de este libro se realizó con
    gbpublisher versión~\gbversion, bajo un modelo de fuente única que
    garantiza la trazabilidad entre el original y cada una de sus
    salidas (HTML, ePub y PDF). Las familias tipográficas utilizadas son
    \gbfuentetexto\ para el texto e \gbfuentetitulos\ para los
    títulos. El proyecto se encuentra disponible en
    \url{https://github.com/albertomoyano/gbpublisher}.
  \end{tcolorbox}}
